% This page is laid out according to Silesian University of Technology's
% requirements for bachelor's thesis
\newpage
\thispagestyle{empty}
\newgeometry{top=2.5cm, bottom=2.5cm, left=3cm, right=2.5cm}
% Measure absolute positioning from margins
\textblockorigin{3cm}{2.5cm}

\begin{onehalfspacing}
\begin{center}
	% Author
    \begin{textblock*}{\textwidth}(0cm, 19cm)
	\begin{flushleft}
	Autor: Maciej Ziaja \linebreak
	Kierujący pracą: dr inż. Sebastian Budzan \linebreak
	\end{flushleft}
	\end{textblock*}
	% Release date and place
	\begin{textblock*}{\textwidth}(0cm, 24cm)
	\fontsize{12}{12} \selectfont
	Gliwice, styczeń 2020
	\end{textblock*}
    % University logo
    \vspace{2\baselineskip}
	\includegraphics[scale=0.75]{polsl_logo_bw_pl}\\
	\vspace{2\baselineskip}
	\fontsize{18}{18} \selectfont
	% University name
	\textbf{\textsc{Politechnika Śląska \linebreak
	Wydział Automatyki, Elektroniki i~Informatyki \linebreak
	Kierunek Automatyka i~Robotyka}} \\
	% Thesis title
	\vspace{3\baselineskip}
	Projekt inżynierski \\
	\vspace{3\baselineskip}
	\fontsize{14}{14} \selectfont
	Wykorzystanie sieci neuronowych do detekcji ziaren
	w~obrazach termowizyjnych
\end{center}
\end{onehalfspacing}
\restoregeometry
